\documentclass[12pt]{article}
\usepackage[T1]{fontenc}
\usepackage[utf8]{inputenc}
\usepackage{lmodern,amsmath,amssymb,amsthm,mathtools}
\usepackage[margin=1in]{geometry}
\usepackage[hidelinks]{hyperref}
\hypersetup{pdftitle={Certified whole-domain bounds and equality barriers for the odd theta–Weil form},
 pdfauthor={Leonard van Hemert},
 pdfsubject={Version 1.0.0; mathematical preprint and certified proof scope}}
\usepackage{microtype}
\setlength{\emergencystretch}{1em}
\newtheorem{theorem}{Theorem}[section]
\newtheorem{proposition}[theorem]{Proposition}
\newtheorem{lemma}[theorem]{Lemma}
\newtheorem{corollary}[theorem]{Corollary}
\theoremstyle{remark}\newtheorem{remark}[theorem]{Remark}
\newcommand{\Ran}{\operatorname{ran}}
\newcommand{\Dom}{\operatorname{Dom}}
\newcommand{\dd}{\,d}
\title{Certified whole-domain bounds and equality barriers\protect\\
for the odd theta--Weil form}
\author{Leonard van Hemert}
\date{13 October 2026 --- Version 1.0.0}
\begin{document}\maketitle
\begin{abstract}
We prove explicit lower bounds on the full complex odd form domain of
a theta-weighted realization of the Riemann--Weil form. After removing
any closed verified kernel containing the score, an admissible residual-enriched trial proves a family of explicit bounds.
One retains $0.132584532412$ of a positive gamma--prime-symbol energy
with loss less than $0.034762$; another retains $0.212458845389$
with loss less than $0.039482$. Among the positive-retention certificates in a set of $231$ tested
parameter pairs, the least scalar loss is $0.031324476155738682727829$, with energy
coefficient $0.000118321042$; no global optimality is asserted. The certificate keeps the
entire physical residual Gram, its coupling outside the trial space,
all prime powers, the contact and pole terms, and the infinite remainders.
The exact obstruction frame has physical rank $56$. We also prove that
the fixed positive midpoint comparison overshoots an actual verified
equality vector by more than $7/5000$ in the original Hilbert norm.
This gives a quantitative barrier to zero-loss certification by that
comparison, even with perfect trial enrichment, and an exact kernel-aware
replacement. The operator algebra uses established residual and
factorization principles; the contribution is their fully priced
theta--arithmetic realization and its equality obstruction. Exact continuum cancellation also gives a complete signed-action
error below $11\cdot10^{-6}$ for the enriched trial, with every prime
tail and pole retained. No zero-error positivity, even-sector conclusion,
or resolution of RH is claimed.
\end{abstract}

\section{The result and its scope}
Put $\Xi(t)=\xi(1/2+it)$, with
$\xi(s)=s(s-1)\pi^{-s/2}\Gamma(s/2)\zeta(s)/2$. For $x\ge0$ define
\[
 \Phi(x)=\sum_{n\ge1}e^{x/2-\pi n^2e^{2x}}
       \bigl(4\pi^2n^4e^{4x}-6\pi n^2e^{2x}\bigr),
 \qquad \Phi(-x)=\Phi(x).
\]
The theta transformation gives $\widehat\Phi=\Xi$ for
$\widehat h(t)=\int_{\mathbb R}h(x)e^{itx}\dd x$.
Set $w^2=\Phi/\cosh(x/2)$, $p=w\sinh(x/2)$ and
$H=H_{\rm odd}\subset L^2(\mathbb R,\dd x)$. The equivalent
positive-half-line measure is $2\dd x$ throughout.
Let $b(t)=\Re\psi(1/4+it/2)-\psi(1/4)$ and let $G$ be the
positive closed form operator $M_wb(D)M_w$. Its form domain is
\[
 \mathcal V=\left\{u\in H:
   \frac1{2\pi}\int b(t)|\widehat{wu}(t)|^2\dd t<\infty\right\}.
\]
The ordinary signed odd form is
\begin{equation}\label{signed}
 \mathfrak a(u)=g(u)+c_\Gamma\|wu\|^2
       -\langle Ju,u\rangle-2|\langle u,p\rangle|^2,
 \quad c_\Gamma=\psi(1/4)-\log\pi,
\end{equation}
where the complete norm-convergent prime-power operator is
\[
 J=\sum_{n\ge2}\frac{\Lambda(n)}{\sqrt n}
       M_w(T_{\log n}+T_{-\log n})M_w.
\]
Here $T_a h(x)=h(x+a)$; both shifts make the convention immaterial
to the selfadjoint sum. The contact constant, gamma symbol, all
prime powers and odd pole in \eqref{signed} are retained. The even
pole disappears only because the present space is odd. The signed realization originates in \cite[Lemma 5.1]{Theta}.
Supplement~S0 proves the domain, kernel and tail facts additionally used here.

Define $q=\Phi'/w$ and $Q=\|q\|^2$. The score is in
$\Dom(G)\cap\ker A_W$, where $A_W$ represents \eqref{signed}.
The positive energy used in this paper is
\begin{align}
 E_8(u)&=\frac1{2\pi}\int d_{8,+}(t)|\widehat{wu}(t)|^2\dd t,
 \label{energy}\\
 d_8(t)&=m_\Gamma(t)-k_8(t),\nonumber\\
 m_\Gamma(t)&=\Re\psi(1/4+it/2)-\log\pi,\nonumber\\
 k_8(t)&=2\sum_{n=2,3,4,5,7,8}\frac{\Lambda(n)}{\sqrt n}
                     \cos(t\log n).\nonumber
\end{align}
This is a different positive form from the score-difference energy
$e_0$ associated with the score. No ordering between them is assumed.

\begin{theorem}[Whole-domain bounds]\label{main}
For every closed verified odd subspace $K\subset\ker A_W$ containing
$q$, every complex $u\in\mathcal V$, and $z=(I-P_K)u$, we prove
\begin{align*}
 \mathfrak a(u)&\ge 0.132584532412 E_8(z)
             -0.034762\|z\|^2,\\
 \mathfrak a(u)&\ge 0.212458845389 E_8(z)
             -0.039482\|z\|^2,\\
 \mathfrak a(u)&\ge\frac{59160521}{500000000000} E_8(z)
             -\frac{31324476155738682727829}{10^{24}}\|z\|^2.
\end{align*}
There are no further finite orthogonality conditions. All constants
use the original physical norm and are independent of the equality-frame
index $m$. Their pointwise maximum is valid.
\end{theorem}
The first two inequalities simultaneously improve the loss and energy
coefficient of the respective earlier bounds
$0.03E_8-0.077614688545235\|\cdot\|^2$ and
$0.18E_8-0.08208813127141\|\cdot\|^2$.
The last bound has the smallest scalar loss among the tested candidates
with strictly positive energy retention,
but its smaller energy coefficient prevents a scalar-only ordering.
A negative lower bound does not establish positivity.
For $z\ne0$, write $r=E_8(z)/\|z\|^2\ge0$. Each assertion is an
ordinary affine lower bound in $r$. At small $r$ the least-loss row is
useful; at larger $r$ a row retaining more energy is stronger. Taking
the maximum is a pointwise operation on these scalar bounds, not the
construction of a new quadratic form. The equality case $z=0$ gives
$\mathfrak a(u)=0$ by the verified kernel identity.

All finite matrices in the construction act on complex coefficient
spaces. Their stored entries happen to be real; a real symmetric
Loewner comparison extends to complex vectors by applying it to their
real and imaginary parts. The contour and form arguments below likewise
use sesquilinear pairings. Thus the complex quantifier is part of the
proof, rather than an inference from tests on real trial vectors.


\subsection{Normalization and the meaning of the quotient}
The norm in Theorem~\ref{main} is
\[
 \|z\|^2=\int_{\mathbb R}|z(x)|^2\dd x
         =2\int_0^\infty|z(x)|^2\dd x.
\]
If the physical Weil test is written as $f=wz$, this norm is
$\int |f(x)|^2/w(x)^2\dd x$. The weight tends rapidly to zero, so
it is essential to state which of these two variables carries the
comparison norm. Multiplication by $w$ also does not impose compact
support. All Fourier cutoffs used in the computation are accompanied
by tails in this original weighted realization.

The word ``verified'' in the kernel hypothesis means that membership
in the operator kernel has been established for that subspace. It
asserts no completeness of the currently known equality vectors.
The smallest permitted choice is $K=\operatorname{span}\{q\}$.
If $K$ is enlarged, $P_K$ remains the actual Hilbert orthogonal
projection; an approximate frame projection cannot be substituted.
The exact signed form annihilates every mixed pairing with $K$, so
\[
 \mathfrak a(u)=\mathfrak a(z),\qquad z=(I-P_K)u.
\]
Moreover $z\in\mathcal V\cap q^\perp$. Thus one certificate on the
score quotient proves the displayed inequality for every such $K$.
This transfer does not assert that the positive energies of two
successive quotients are ordered, and it does not transfer a numerical
Gram or conditioning bound onto a newly projected frame.

There are three distinct spaces in the finite construction. The
coefficient space of $B$ has dimension $2906$; the exact map $C=BZ$
has $56$ input coordinates; and its vectors lie in the infinite
physical Hilbert space. The coefficient metric $N=Z^*Z$ measures
$\|Zx\|^2$, whereas $H=C^*C$ measures $\|Cx\|^2$. The proof retains
both. This distinction is particularly important because the physical
frame has a much smaller lower Gram bound than its coefficient frame.
Numerical orthogonality in one metric cannot supply an inverse or a
form comparison in the other.


\section{A complete matrix residual certificate}
The next lemma records the precise reusable certificate used below.
It is an operator-valued version of standard coercive residual
certification; see \cite[Section 4.3, Proposition 4C]{PR}.
\begin{lemma}\label{matrixlemma}
Let $E\ge0$ be selfadjoint on a Hilbert space, $s>0$, and $C$ a
finite column map. Let $N>0$ be a coefficient matrix and $\nu>0$.
For a trial map $Y_0$ with columns in $\Dom(E)$, suppose a
selfadjoint proxy $E_m$ is defined on its range and
$\|(E-E_m)|_{\Ran Y_0}\|\le\eta$. Set
\[
 Q_0=Y_0^*Y_0,\quad r=C-(E_m+s)Y_0,\quad S_0=r^*r,
 \quad K_m=C^*Y_0+Y_0^*C-Y_0^*(E_m+s)Y_0.
\]
For any $\varepsilon>0$, the matrix comparison
\begin{equation}\label{matrixcriterion}
 \nu\left[K_m+
   \left(\eta+\frac{(1+1/\varepsilon)\eta^2}s\right)Q_0
    +\frac{1+\varepsilon}sS_0\right]\le\theta N
\end{equation}
implies $\nu CN^{-1}C^*\le\theta(E+s)$ as forms. It retains the
entire residual in the ambient Hilbert space.
If symmetric centers $K_c,Q_c,S_c$ have full operator errors
$k_e,q_e,s_e$, adding
$\nu[k_e+(\eta+(1+1/\varepsilon)\eta^2/s)q_e
 +(1+\varepsilon)s_e/s]I$ to the center matrix suffices.
\end{lemma}
\begin{proof}
Put $T=E+sI$. Selfadjointness and $T\ge sI$ give a bounded
inverse with $0\le T^{-1}\le s^{-1}I$. This is precisely where the
strictly positive shift is used. Since each trial column belongs to
$\Dom(T)$, the maps $TY_0$ and $R=C-TY_0$ are well defined in the
ambient Hilbert space. Expanding the entire residual gives
\begin{align*}
 R^*T^{-1}R
 &= (C-TY_0)^*T^{-1}(C-TY_0)\\
 &= C^*T^{-1}C-C^*Y_0-Y_0^*C+Y_0^*TY_0.
\end{align*}
Hence, with $K=C^*Y_0+Y_0^*C-Y_0^*TY_0$,
\begin{equation}\label{completed}
 C^*T^{-1}C=K+R^*T^{-1}R\le K+s^{-1}R^*R.
\end{equation}
There is no Galerkin equation in this identity. An arbitrary admissible
trial works, and its quality affects only the residual size.

Let $D=(E-E_m)Y_0$. For every coefficient vector $x$ the assumed
action bound gives $\|Dx\|\le\eta\|Y_0x\|$. Since the quadratic
pairing is real by symmetry,
\[
 K-K_m=-\operatorname{Herm}(Y_0^*D)\le\eta Q_0,
 \qquad D^*D\le\eta^2Q_0.
\]
Also $R=r-D$. Applying scalar Young's inequality to the cross term
in $\|rx-Dx\|^2$ for every $x$ proves the matrix inequality
\[
 R^*R\le(1+\varepsilon)S_0
                   +(1+1/\varepsilon)\eta^2Q_0.
\]
Substitution in \eqref{completed} proves
$\nu C^*T^{-1}C\le\theta N$ from \eqref{matrixcriterion}.

For completeness, the passage from this coefficient inequality to
a form inequality is exact. Define the bounded map
\[
 L=\sqrt\nu\,T^{-1/2}CN^{-1/2}.
\]
The coefficient statement is $L^*L\le\theta I$, equivalently
$\|L\|^2\le\theta$, and therefore $LL^*\le\theta I$ on the full
Hilbert space. For an arbitrary $z\in\Dom(T^{1/2})$, test the latter
inequality on $T^{1/2}z$. Since $T^{-1/2}T^{1/2}z=z$, this yields
\begin{align*}
 \nu\langle CN^{-1}C^*z,z\rangle
 &=\|L^*T^{1/2}z\|^2\\
 &\le\theta\|T^{1/2}z\|^2.
\end{align*}
This is the claimed inequality on the entire form domain. It is an
instance of standard majorization and factorization
\cite[Theorem 1]{Douglas}, with its hypotheses supplied by the shift
and $N>0$. The physical matrix $C^*C$ never replaces $N$ in this step.

Finally, if $K_m\le K_c+k_eI$, $Q_0\le Q_c+q_eI$ and
$S_0\le S_c+s_eI$, all their multipliers in
\eqref{matrixcriterion} are positive. Substitution therefore adds
exactly the stated scalar error times $I$. The center matrices need
not themselves be positive for this upper-bound argument.
\end{proof}

Taking separate largest eigenvalues of $K_m$ and $S_0$ before adding
them loses their correlation. Formula \eqref{matrixcriterion} retains
that correlation until the final finite comparison. The difference is
decisive in the actual calculation: the separate-norm certificate fails
at $\mu=1/20$, while the full matrix certificate passes there.


\subsection{From a contraction to an energy--loss bound}
Suppose on a closed quotient the original form satisfies
\[
 \mathfrak a(z)\ge E(z)-\langle Rz,z\rangle-\rho\|z\|^2,
 \qquad R=\nu CN^{-1}C^*.
\]
If the preceding certificate gives $R\le\theta(E+sI)$, with
$0\le\theta<1$, substitution retains the positive energy:
\begin{align}\label{floorconversion}
 \mathfrak a(z)
 &\ge E(z)-\theta\bigl(E(z)+s\|z\|^2\bigr)-\rho\|z\|^2\\
 &= (1-\theta)E(z)-(\rho+\theta s)\|z\|^2.\nonumber
\end{align}
The parameter $s$ is a shift of the positive operator. It is not a
positive gap for the signed operator. Writing $s=\mu-\rho$ is a
convenient parametrization with $\mu>\rho$; it gives the loss
$\rho+\theta(\mu-\rho)$. A successful comparison improves the
scalar floor relative to $-\mu$, while keeping $1-\theta$ of $E$.
If instead $\theta\ge1$, the same algebra is valid but its energy
coefficient is nonpositive. This explains the positive-retention
restriction in the least-loss statement.

The construction has two logically separate tasks. First, an analytic
and interval calculation supplies the finite-rank price $R$ and the
whole remainder $\rho$. Second, a residual comparison absorbs most
of $R$ into a shifted positive energy. A small finite signed eigenvalue
alone performs neither task: it does not control the physical coupling
to the rest of the Hilbert space. Both tasks are carried out below.

\section{The actual theta--arithmetic certificate}
Put $r_8=(-d_8)_+$ and $R_8=M_wr_8(D)M_w+2p\otimes p$.
The literal common signed form is
\begin{equation}\label{split}
 \mathfrak a(u)=E_8(u)-\langle R_8u,u\rangle
                         -\langle(J-J_8)u,u\rangle.
\end{equation}
Since $|d_{8,+}-b|\le|c_\Gamma|+\|k_8\|_\infty$, $E_8-G$ is
bounded, $\Dom(E_8)=\Dom(G)$ and $\Dom(E_8^{1/2})=\mathcal V$.
Let $S=I-q\otimes q/Q$ and $E_Y$ be its positive form compression
to $Y=q^\perp$. The projection preserves the graph domain.

The exact finite data consist of $2905$ nonzero frequency cells
$I_i=[k_i/64,(k_i+1)/64]$, their midpoints $t_i=(2k_i+1)/128$,
and the exact upward binary64 mass bounds $\mu_i^+$ for
$\int_{I_i}r_8(t)\dd t/\pi$. Define
\begin{gather*}
 Be_i=\sqrt{\mu_i^+}S[w\sin(t_ix)],\qquad Be_p=\sqrt2Sp,\\
 C=BZ,\qquad N=Z^*Z,\qquad R_{56}=\frac{11}{10}CN^{-1}C^*.
\end{gather*}
Every entry of the supplied $2906\times56$ matrix $Z$ is interpreted
as an exact dyadic rational. Its origin in a numerical solve is not
an orthogonality or rank premise.

\begin{proposition}[Verified realization data]\label{inputs}
The supplied interval certificates and analytical estimates establish
\begin{gather*}
 (1-10^{-8})I\le N\le(1+10^{-8})I,\qquad
 C^*C\ge\frac{21}{10^8}N,\quad C^*C<\frac{32}{125}I,\\
 \mathfrak a(z)\ge E_8(z)-\langle R_{56}z,z\rangle-\rho\|z\|^2,
 \quad \rho=\frac{979236349}{2000000000000},\quad z\in\mathcal V\cap Y.
\end{gather*}
The actual physical rank is $56$. With all $2469$ positive zeros
through $3000$, both signs and multiplicities included, the positive
sampling operator $S_{3000}$ satisfies
\[
 \|(A_W-S_{3000})|_{\Ran C}\|<1.4\cdot10^{-18},\qquad
 \mathfrak a|_{\Ran C}\ge\frac{38}{10^{16}}I.
\]
For $R_m=BB^*$ and $E_m=S_{3000}+R_m$, the whole action price on
this trial range is
\[
 \|(E_Y-E_m)|_{\Ran C}\|<\eta,
 \qquad \eta=0.013155646076930.
\]
\end{proposition}
\begin{proof}
We give the complete bridge and the finite enclosures which make the
data reviewable. Supplement~S1--S7 and the certificate files identified in the
manifest give the interval construction, including all infinite
allowances; none is a sampled-matrix positivity assumption.

The full support certificate proves $r_8(t)=0$ for $|t|\ge10000$,
its positive-half-line mass is below $14$, and
$\|J-J_8\|\le t_8=32237/10^{12}$. A full theta moment gives
\[
 \int\|S[w\sin(tx)]-S[w\sin(t_ix)]\|^2r_8(t)\dd t/\pi
 \le\delta_0=21/1024000.
\]
Young's inequality on every cell, with parameter $1/10$, gives
$SR_8S\le(11/10)BB^*+11\delta_0I$, including the identical pole.
The positive coefficient upper Gram $\Gamma$ retains every
frequency difference, sum, score and pole correlation. The full
dense residual certificate proves
$\|\Gamma-Z\Lambda Z^*\|<3/12500$, with $\Lambda>0$ exact dyadic.
Compression to $(\Ran Z)^\perp$ therefore gives
$\|B(I-P)B^*\|\le3/12500$, where $P=ZN^{-1}Z^*$.
Adding the full prime tail yields
$\rho=t_8+11\delta_0+(11/10)(3/12500)$, exactly as displayed.

The actual feature Gram is
$B^*B=D_\mu(G_0-b_0b_0^*/Q)D_\mu$, with
$G_0(i,j)=[\widehat{w^2}(t_i-t_j)-\widehat{w^2}(t_i+t_j)]/2$,
$G_0(p,i)=2\int_0^\infty w^2\sinh(x/2)\sin(t_ix)\dd x$,
$G_0(p,p)=\tau$, $b_{0,i}=-t_i\Xi(t_i)$, $b_{0,p}=-1/4$.
New outward integrals use the actual $Q\approx11.06766574714$
and $\tau\approx0.005712908530$, with complete allowances
$2^{-250}+2^{-180}$. The whole compression error is below
$3.2\cdot10^{-9}$. A binary spectral locator is corrected by its full
orthogonality defect and residual, giving the rational Gram bounds
above. Thus $\kappa_B\le\sqrt{10^8/21}<2182.178903$.

Native FLINT $3.6.0$ zero enclosures exhaust its multiplicity count
$N(3000)=2469$ \cite{FLINT}. Disjoint actual roots prove finite-range
simplicity. With $V_0e_\gamma=w\sin(\gamma x)$ and $A_0=V_0^*C$,
the complete paired moments are $L=2A_0^*A_0$ and
$O=4A_0^*(V_0^*V_0)A_0$. All $2469^2$ correlations are kept.
A new order-$20$, spacing-$1/8$ Taylor grid through $128$ uses
$320$-node outward quadrature, the entire lattice and physical tails.
The complete bounds are $\|\Delta A_0\|_F\le2\cdot10^{-12}$,
$\|\Delta(V_0^*V_0)\|\le3.2\cdot10^{-11}$,
$\|\Delta L\|\le10^{-11}$ and $\|\Delta O\|\le3.1\cdot10^{-10}$.
Direct rectangular-factor certification proves $L\ge10^{-15}I$.
The maximal actual midpoint is $T=313317/128$. The full contour
and zero-count tails at height $L_0=3000$ are
\begin{align*}
 e&=10^6\kappa_B^2e^{T/4}Z_{1/4}(L_0),\\
 d&=11000\kappa_Be^{T/8}Z_{1/8}(L_0),\\
 Z_a(L)&=e^{-aL}\left[2L\log L+
                \frac{2(\log L+1)}a+\frac4{a^2L}\right].
\end{align*}
They give $e<3\cdot10^{-43}$ and $d<1.4\cdot10^{-18}$, including
every omitted/off-line zero and multiplicity. This proves the finite
signed gap without asserting any infinite-block inverse.

Finally introduce the true-mass midpoint map on the same feature input.
Its distance from the continuum feature map is at most $\sqrt{\delta_0}$.
Thus its price differs from $SR_8S$ in norm by at most
$2\sqrt{0.255985513\delta_0}+\delta_0$.
The entire weighted mass-deficit certificate adds less than
$0.00855266$, equivalently $427633/50000000$,
with zero pole deficit.
On $Y$, $E_Y=A_W+SR_8S+S(J-J_8)S$.
Adding $t_8$ and the whole action tail $d$ proves the stated $\eta$.
The columns are in $\Dom(G)$ by complete theta smoothing and
verified score projection. This proves the domain and action premises.
\end{proof}

\subsection{The coupled finite comparison}
For each shift $\mu>\rho$, put $s=\mu-\rho$ and choose the supplied
exact binary trial $Y_0=CX$. It needs no certified solve accuracy.
Let $G_B=B^*B$, $G_z=V_0^*V_0$, $Q_b=V_0^*B$. The whole residual is
\begin{align*}
 r&=BD_1+V_0D_2,\\
 D_1&=Z(I-sX)-G_BZX,\qquad D_2=-2A_0X,
\end{align*}
with exact Gram
\begin{align*}
 r^*r={}&D_1^*G_BD_1+D_1^*Q_b^*D_2\\
       &+D_2^*Q_bD_1+D_2^*G_zD_2.
\end{align*}
All $2469\cdot2906$ mixed correlations appear. In particular the
component outside $\Ran C$ is never removed. The certificate prices
every input, inner/outer product, exact shift conversion and coefficient
assembly. If $e_d$ bounds the physical coefficient-map error and
$r_e$ the center Gram error, its actual Gram upper price is
$r_e+2\sqrt{\lambda_{\max}(S_c)+r_e}\,e_d+e_d^2$.
This price is included before Lemma~\ref{matrixlemma} is used.
All small matrix locators use exact-input $192$-bit Arb residuals and
their polar-unitary correction. Only locators come from floating point.

The complete upper prices are exact rationals:
\begin{center}
\begin{tabular}{c|c|c|c}
$\mu$ & $\varepsilon$ & matrix price $p$ & $\theta$\\\hline
$1/20$ & $2/3$ & $100157725018/125000000000$ & $400631/500000$\\
$1/25$ & $1/2$ & $450253581783/500000000000$ & $225127/250000$
\end{tabular}
\end{center}
Each $p$ includes the full finite error prices and an additional
$10^{-6}$ reserve. The strict margins to $\theta(1-10^{-8})$ are
$9592169/(5\cdot10^{13})$ and $20685723/(25\cdot10^{12})$.
Thus Lemma~\ref{matrixlemma} proves $R_{56}\le\theta(E_Y+s)$.
Inserting this into Proposition~\ref{inputs} gives
$\mathfrak a\ge(1-\theta)E_8-[\rho+\theta(\mu-\rho)]\|\cdot\|^2$,
which proves the two original matrix-only bounds on $q^\perp$.
The enriched improvements in Theorem~\ref{main} follow in the next section.
For a larger verified $K$, $z\in\mathcal V\cap Y$ and the literal
kernel identity gives $\mathfrak a(u)=\mathfrak a(z)$.
This transfers the full theorem without recertifying a changed frame.

\section{Admissible residual enrichment and its full action}
The original frame-only action error cannot be applied to a new
residual direction. We instead use the actual dyadic trial
$Y=Bb+V_0v$ with $56$ input columns, selected from the residual outside
$\Ran C$. No inverse or rank of an enlarged trial Gram is assumed.
For every fixed frequency $a$, $w^2\sin(ax)$ is Schwartz; the complete
theta series proves the same smoothing for the pole and score terms.
The bound $b(t)\le17+9\log(1+|t|)$ therefore places each column in
$\Dom(G)\cap q^\perp$. The existing positive physical Gram justifies the
projection off $\Ran C$ used in selecting the residual. Exact binary
selected coefficients themselves define the trial, so rounded numerical
orthogonality is not a premise.

Put $\ell_8=m_\Gamma-k_8$, $L_8=M_w\ell_8(D)M_w$, and let
$\mathcal E$ denote the positive compressed operator previously called
$E_Y$. Its full domain identity is
\[
 \mathcal E=SL_8S+R_m-2Sp\otimes Sp+\Delta_R,
 \qquad \|\Delta_R\|\le\epsilon_R,
\]
where
$\epsilon_R=2\sqrt{0.255985513\delta_0}+\delta_0+0.00855266$.
This follows from $d_{8,+}=\ell_8+r_8$ and the already proved
whole-space mass and midpoint prices. It makes no low-zero action
assumption on the new columns.

The completed outward Fourier calculation supplies a whole map
$\widehat A$ with
\[
 \|\mathcal E Y-\widehat A-\Delta_RY\|\le g,
 \quad g<1.421821\cdot10^{-5},\qquad \|Y\|<2.159290.
\]
Here the finite lattice has spacing $1/20$, frequency height $3300$,
and physical radius $3$. The theta strip envelope bounds the complete
Fourier row integral by $1.001591\cdot10^9$; contour shifting prices
the alias below $1.532843\cdot10^{-5}$ before $\|w\|$.
Complete node/location/rounding price is below
$1.797275\cdot10^{-8}$, omitted frequencies below $10^{-88}$,
and physical tails below $10^{-157}$.
The physical action Gram retains the full $\widehat{w^2}(t-s)$
Toeplitz convolution, both Fourier shifts and the score subtraction.
Native outward complex DFT radii are included. Supplement~S9 gives
all analytic envelopes and product guards; S10 gives the whole-action calculation and its error ledger. These are operator-map
bounds, not sampled form positivity.

\subsection{The entire action error in the resolvent comparison}
For $s=\mu-\rho>0$, set
\[
 G_Y=Y^*Y,\quad CY=C^*Y,\quad YA=Y^*\widehat A,
\]
\[
 CA=C^*\widehat A,\qquad AG=\widehat A^*\widehat A,
 \qquad H=C^*C.
\]
These are exact physical matrices. Their outward numerical enclosures
will be inserted only after the operator calculation. Define
\begin{align*}
 \widehat R_s&=C-\widehat A-sY,\\
 K_s&=CY+CY^*-sG_Y-(YA+YA^*)/2,\\
 S_s&=\widehat R_s^*\widehat R_s\\
    &=H-CA-CA^*-s(CY+CY^*)+AG\\
    &\hspace{12mm}+s(YA+YA^*)+s^2G_Y.
\end{align*}
The appearance of $AG$ is essential: $\widehat A$ takes values in
the full Hilbert space, so this product measures action outside the
trial span as well as inside it. Replacing $AG$ by a compressed action
product would no longer equal the displayed residual Gram.

Write the exact action discrepancy as
\[
 \mathcal EY=\widehat A+\Delta_RY+D,
 \qquad\Delta_R^*=\Delta_R,\quad
 \|\Delta_R\|\le\epsilon_R,\quad\|D\|\le g.
\]
The two errors have different structure. The first is a bounded
selfadjoint operator applied to $Y$, so its quadratic price keeps
$G_Y$. The second is a complete map error on the coefficient input,
so its worst quadratic price is a scalar multiple of $I$.
For each coefficient vector $x$,
\begin{align*}
 |\langle\Delta_RYx,Yx\rangle|
     &\le\epsilon_R\|Yx\|^2,\\
 |\Re\langle Dx,Yx\rangle|
     &\le g\|Y\|_{\rm upper}\|x\|^2.
\end{align*}
Consequently the exact completed-square quadratic is bounded above by
$K_s+\epsilon_RG_Y+g\|Y\|_{\rm upper}I$.

For any $\tau>0$, vector Young's inequality also gives
\begin{align*}
 (\Delta_RY+D)^*(\Delta_RY+D)
 \le (1+\tau)\epsilon_R^2G_Y+(1+1/\tau)g^2I.
\end{align*}
The true residual is $\widehat R_s-(\Delta_RY+D)$. Apply Young a
second time, now with $\beta>0$, to separate this entire residual
from its action error. Combining with \eqref{completed} gives
\begin{align}\label{re5main}
 C^*(\mathcal E+s)^{-1}C\le{}&K_s+a_YG_Y+a_SS_s+a_II,\\
 a_Y&=\epsilon_R+
          \frac{(1+1/\beta)(1+\tau)\epsilon_R^2}{s},\nonumber\\
 a_S&=\frac{1+\beta}{s},\nonumber\\
 a_I&=g\|Y\|_{\rm upper}
       +\frac{(1+1/\beta)(1+1/\tau)g^2}{s}.\nonumber
\end{align}
The selected value $\tau=g/(\epsilon_R\|Y\|_{\rm upper})$ is
strictly positive. No rank or inverse of $G_Y$ occurs anywhere in
this proof. The coefficient space of $Y$ has the same dimension as
that of $C$, which is all that is required by these products.

\subsection{Inserting finite matrix uncertainties}
For clarity, let the exact matrices above have symmetric or Hermitian
centers carrying operator errors
$e_G,e_{CY},e_{YA},e_{CA},e_{AG},e_H$.
With exact rational $s$, the triangle inequality gives the input-error
parts
\begin{align*}
 e_K&\le2e_{CY}+se_G+e_{YA},\\
 e_S&\le e_H+2e_{CA}+2se_{CY}+e_{AG}
                  +2se_{YA}+s^2e_G.
\end{align*}
The implementation additionally encloses the products and sums that
form those centers, including conversion of the shift. Such assembly
errors are added to $e_K,e_S$; they are not absorbed into the map
error $g$ without proof. If $U_c$ is the center of
$K_s+a_YG_Y+a_SS_s$, then a sufficient final upper comparison is
\[
 \frac{11}{10}\left[U_c+
       (e_K+a_Ye_G+a_Se_S+a_I)I\right]\le\theta N.
\]
This follows because $a_Y,a_S\ge0$. An exact-input residual and
orthogonality correction encloses the largest eigenvalue of the
center. Its certified upper endpoint, after all scalar allowances,
is compared with $\theta(1-10^{-8})$, using Proposition~\ref{inputs}.
Thus a numerical spectral locator alone is never the final premise.
The full input allowances and product guards are given in
Supplement~S10--S11.

\subsection{The certified parameter family}
Supplement~S11 certifies the endpoints for $231$ rational $(\mu,\beta)$ candidates from
these same cached Grams with independent exact-input Arb residuals,
polar corrections, full assembly prices and rational endpoints.
For $\mu=1/25,\beta=3/2$, the certified $\theta$ is exactly
$0.867415467588$; for $\mu=1/20,\beta=1277/512$, it is exactly
$0.787541154611$. Their exact losses are below the upward decimals in
Theorem~\ref{main}. For $\mu=401/12800,\beta=1913/2048$,
$\theta=499940839479/500000000000$ gives the third exact bound.
Each final price is strictly below $\theta(1-10^{-8})$.
Factorization and the original rest price prove Theorem~\ref{main}.
The accompanying finite verifier checks the exact candidate endpoints;
the saved primitive Grams determine each selected comparison.

The central upper matrix has the form $A+\beta sG_Y+M/s$, $M\ge0$.
It is operator convex in $s$, but its derivative can have either sign.
Thus the adaptive rational-neighborhood search makes no monotone
bisection or interval exclusion claim. Only candidates with $\theta<1$ enter the stated least-loss comparison;
other tested endpoints can have a negative energy coefficient. Best tested
in this class is not proved optimal.
At the least-loss candidate the central price is about $0.997674$;
the complete finite uncertainty is below $0.002208$. The actual whole
residual and bounded continuum error remain substantive prices.

\section{An actual equality obstruction to the fixed positive price}
Define the first verified pole-null jet $h=(\Phi'''-\Phi'/4)/w$,
$a=\langle h,q\rangle/Q$, and $k=h-aq$. Both $h$ and $q$ belong
to $\Dom(G)\cap\ker A_W$; their transforms are polynomial multiples
of $\Xi$, vanishing at every actual zero, with the full signed
explicit formula and graph-domain extension. Moreover
$\widehat{wk}(t)=it(t^2+1/4+a)\Xi(t)$, so $k\ne0$.
\begin{theorem}[Fixed-price equality barrier]\label{barrier}
On the score-only quotient, the actual vector $k$ satisfies
\[
 \langle R_{56}k,k\rangle-E_Y(k)>\frac7{5000}\|k\|^2,
 \qquad \mathfrak a(k)=0.
\]
Therefore any comparison $R_{56}\le\theta(E_Y+s)$ with
$0<\theta\le1$ requires $\theta s>7/5000$. Its scalar price
$\alpha=\rho+\theta s$ must exceed $0.0018896181745$.
This is a barrier for this fixed positive comparison, not a negative
direction of the actual Weil form or a barrier after removing $k$.
\end{theorem}
\begin{proof}
Two new complete meromorphic Arb integrals enclose
$\langle h,q\rangle\approx73.93892065366$ and
$\|h\|^2\approx54152.69071824$; their full theta and physical
allowances are less than $2^{-130}$. They give $a>6$ and
$0<D=\|k\|^2<55000$. The degree-five theta-atom jet polynomial is
$-32y^5+360y^4-1056y^3+810y^2-90y$. Its complete lattice comparison
bounds $|wh|/\Phi<20000e^{6x}$. The entire doubled norm and cross
tails beyond $x=2$ are below
$10^{11}\int_{54}^\infty t^8e^{-3t}\dd t<2^{-140}$.
On $[0,2]$ the omitted lattice and quotient expansion cost less than
$2^{-170}$; both costs are added before projection.

Every midpoint price term is nonnegative. Keep only the $635$ cells
with $t_i\le20$, replace $a$ by $6$ in the squared sine pairing,
and evaluate the exact upward dyadic masses and actual $\Xi(t_i)$
with complex Arb arithmetic. The resulting sum is greater than $1875$.
Hence $r=\langle R_mk,k\rangle/D>3/88$; the omitted pole and all
other midpoint cells only increase this lower bound.
To compare this midpoint energy with the actual continuum energy,
let $F_0$ be the continuum feature map and $F_1$ the midpoint map
using the true cell masses, on a common $L^2$ coefficient space.
Their positive products are respectively $SR_8S$ and $R_t$, and
$\|F_0-F_1\|\le\sqrt{\delta_0}$. Since $0\le R_t\le R_m$,
\begin{align*}
 \sqrt{\langle SR_8Sk,k\rangle}
 &=\|F_0^*k\|\\
 &\le\|F_1^*k\|+\sqrt{\delta_0}\|k\|\\
 &\le\sqrt{\langle R_mk,k\rangle}+\sqrt{\delta_0}\|k\|.
\end{align*}
This vectorwise bound uses no sign assertion about the difference
$SR_8S-R_m$. It is stronger for the present purpose than replacing
that difference by a whole-space absolute norm.

Put $\varepsilon_c=3/12500$ and $T=S(J-J_8)S$. The exact identities
and inequalities on the kernel vector are
\begin{align*}
 E_Y(k)&=\langle SR_8Sk,k\rangle+\langle Tk,k\rangle,\\
 \langle Tk,k\rangle&\le t_8D,\\
 \langle R_{56}k,k\rangle
 &\ge\frac{11}{10}\langle R_mk,k\rangle
                  -\frac{11}{10}\varepsilon_cD.
\end{align*}
The first follows from the signed form and $\mathfrak a(k)=0$;
the last is the complete dense capture. Dividing by $D>0$ and
using the feature-map triangle bound proves
\begin{align}\label{barrierfunction}
 \frac{\langle R_{56}k,k\rangle-E_Y(k)}D&\ge f(r),\\
 f(r)&=\frac r{10}-2\sqrt{\delta_0r}-\delta_0
                     -\frac{11}{10}\varepsilon_c-t_8.\nonumber
\end{align}
For $r\ge3/100$, its derivative
$1/10-\sqrt{\delta_0/r}$ is positive. The outward scalar comparison
$f(3/88)>7/5000$, together with $r>3/88$, gives the first assertion.
Every signed arithmetic term is present in the kernel identity; only
nonnegative midpoint terms were omitted in obtaining the lower sum.

Finally suppose $R_{56}\le\theta(E_Y+s)$ with $0<\theta\le1$.
Evaluation on this same vector, subtraction of $E_Y(k)$, and division
by $D$ give
\[
 \frac7{5000}<\theta s-(1-\theta)\frac{E_Y(k)}D
                                    \le\theta s.
\]
The nonnegativity of $E_Y$ justifies the last inequality. Adding the
fixed $\rho$ proves the strict scalar-loss barrier in the theorem.
Supplement~S8 gives the scalar integral enclosures and complete tails.
\end{proof}

The barrier identifies lost equality cancellation caused by the positive
midpoint enlargement. It redirects improvement toward the signed rest
or a kernel-aware price, rather than unlimited decimal tuning.
For $\Pi=I-k\otimes k/D$, set $B'=\Pi B$, $C'=\Pi C$.
The exact update is
\begin{align*}
 B'^*B'&=B^*B-b_kb_k^*/D,\\
 C'^*C'&=C^*C-c_kc_k^*/D,\qquad b_k=B^*k,\quad c_k=Z^*b_k.
\end{align*}
Actual-zero orthogonality gives $V_0^*k=0$; hence the sampling,
zero Gram and zero/feature cross-Gram remain unchanged. The projected
rank is still $56$, since the existing sampling factor is injective.
All columns retain their graph domain. The full projected action error
costs $d\|CX\|+\eta_c\|C'X\|$; replacing the first norm by the
second without proof would be invalid. Supplement~S8 proves these kernel-aware identities. Its first complete numerical
test has upper price greater than one and proves no improved sign.

\section{Signed continuation and the remaining complement}
The positive comparison above is sufficient for Theorem~\ref{main}, but
its equality barrier suggests retaining the signed cancellation. On the
score quotient the exact operator-domain identity is
\[
 A_Y=SL_8S-2Sp\otimes Sp-S(J-J_8)S.
\]
Thus the continuum midpoint error cancels before approximation. For the
same admissible enriched trial, Supplement~S12 proves a complete signed
action error below $11/10^6$, including the full prime tail and physical
output. Supplement~S13 gives a different cached signed approximant with
error below $7/12500$, an actual positive five-dimensional subspace with
form lower bound $\|f\|^2/5000$, and its full outside coupling. These
secondary conclusions do not establish the sign of the complement.
Supplement~S14 states an inverse-free, fully coupled Schur criterion and
identifies its unverified mixed-action hypotheses. No improvement beyond
Theorem~\ref{main} is inferred from them.

\section{Reproducibility, literature and what remains}
The accompanying release includes exact binary trials and matrices,
input and executed-source fingerprints, the analytical supplement and a
bounded finite verifier. That verifier checks hashes, rational consequences,
selected enriched comparisons reconstructed from cached matrices, and the
robust-capture inverse/Rayleigh calculation. It accepts the upstream
quadrature, zero-completeness, Taylor, large-Gram and whole-action error
allowances as inputs; it does not regenerate those computations.
The distinction between analytical proof, certificate inputs, finite replay
and formal verification is recorded explicitly in the release verification report.
Ordinary sampled positivity and an unpriced FFT are not certificate inputs.

The cancellation of a common bounded operator is elementary; its
quantitative benefit here depends on the actual whole-action ledger.
The factorization, coercive residual square and legal Schur inverse are
standard \cite{Douglas,PR,DSS}. Our specific result concerns the whole
theta-weighted complex odd domain and the exact arithmetic realization.
Zhu's compact-window theorem already certifies finite/tail coupling and
odd-sector positivity \cite[Theorems 1.1--1.2 and Section 6]{Zhu}.
Specifically, Zhu's Theorem~1.1 lifts a finite Legendre block with
explicit tail and coupling costs; Theorem~1.2 and Section~6 certify
positivity for complex tests supported in $[-0.8,0.8]$. The odd-sector
floor there is $8.2065\cdot10^{-15}$. Here the test is $wu$ on the whole
line, the comparison norm is $\|u\|_{L^2(dx)}$, and the arithmetic tail
is infinite. Compact support and theta multiplication impose different
conditions, so neither numerical constant dominates the other theorem.
Suzuki's Friedrichs realization on $L^2(-a,a)$ uses the same gamma
multiplier, but its small-window positivity theorem has different
support and norm hypotheses \cite[Theorems 1.1 and 1.4]{Suzuki}.
Connes--Consani's archimedean trace lower bound has a finite loss under
compact multiplicative support and a Fourier vanishing condition; it
does not include the present full prime operator
\cite[Theorem 6.11]{CC}. These are structural precedents. The contribution
claimed here is the explicit theta-weighted whole-domain certificate
and the fixed-price equality obstruction. No first-method or exhaustive
priority claim is made.

The zero-error task remains the original signed comparison
$E_8(z)\ge\langle(R_8+J-J_8)z,z\rangle$ on the required quotient.
At fixed $\rho>0$, even $s\downarrow0$ in a successful sufficient
comparison would retain the floor $-\rho$; Theorem~\ref{barrier}
locates a stronger obstruction for this particular fixed price.
The residual-enriched whole-action certificate is used in
Theorem~\ref{main}. The signed continuation certifies an actual positive
five-dimensional subspace and its whole outside coupling; neither it
nor the sharper signed-action estimate proves the required complement sign. The even sector and all-index Li transfer remain
necessary. The published metric congruence retains its full
$\beta\|M_\beta^{-1/2}v\|^2$ term and fixed-frame logarithmic
conditioning \cite[Proposition 7.1]{Theta}; no $\beta\to0$ control follows.

\subsection{Formal verification scope}
The accompanying executable Lean~4.32.0 project checks the central
abstract argument, with mathlib pinned to a recorded commit. It proves
the coercive inverse bound, exact full-residual identity and upper bound,
the complete action-error comparison, and the coefficient-metric
transfer on the entire form domain. The composed declaration
\texttt{enriched\_certified\_retained\_energy} passes from the full action
decomposition and its finite upper-price certificate to the
energy--loss conclusion. Complex scalar restriction retains every
complex vector and its Hermitian quadratic; it does not replace the
coefficient metric by the physical Gram.

The kernel also checks the exact rational endpoint identities and pass
predicates for all $231$ saved records, and the three displayed numerical
consequences under their explicit premises. The $22$ audited declarations
have only \texttt{propext}, \texttt{Classical.choice} and
\texttt{Quot.sound} in their actual transitive axiom reports, with no
\texttt{sorryAx}, custom axiom or native-evaluation extension.

This is conditional formal verification of the proof mechanism and exact
arithmetic. Construction of the theta operators and closed forms,
resolvent existence, admissibility of the actual trial columns,
theta/Fourier identities, prime and zero tails, interval matrix enclosures,
and original-form/kernel identification remain analytical inputs proved
or certified in the written account. The entire theta--Weil specialization
and the equality-barrier proof are not claimed as formalized.
The included \texttt{lean/README.md} and
\texttt{verification/FORMAL\_CORRESPONDENCE.md} map the declarations to the
published statements; the build and actual axiom output accompany the sources.

\begin{thebibliography}{9}
\bibitem{Theta} L. van Hemert, \emph{Sharp logarithmic coupling estimates
for theta equality frames in a signed Weil-form realization}, v1.0.0,
11 October 2026, \href{https://doi.org/10.5281/zenodo.23298499}
{doi:10.5281/zenodo.23298499}, Lemma 5.1 and Proposition 7.1.
\bibitem{Douglas} R. G. Douglas, \emph{On majorization, factorization,
and range inclusion of operators on Hilbert space}, Proc. AMS 17 (1966),
413--415, Theorem 1. \href{https://home.agh.edu.pl/~rudol/Operat/DouglasFactorizationLemma.pdf}{Original scan}.
\bibitem{PR} A. T. Patera and G. Rozza, \emph{Reduced Basis Approximation
and A Posteriori Error Estimation for Parametrized Partial Differential
Equations}, Version 1.0, Copyright MIT 2006--2007, Part I,
Section 4.3, Proposition 4C and (4.26b).
\href{https://people.sissa.it/~grozza/wp-content/uploads/2024/10/PR07_Patera_Rozza_book_part1_redux.pdf}{Author-hosted text}.
\bibitem{DSS} G. Dusson, I. M. Sigal and B. Stamm,
\emph{The Feshbach--Schur map and perturbation theory},
\href{https://arxiv.org/abs/2105.02058v1}{arXiv:2105.02058v1}, Theorem 1.2.
\bibitem{Zhu} X. Zhu, \emph{Weil positivity in compact windows: a finite
reduction, certified two-sided bounds, and a Landau--Widom decay law},
\href{https://arxiv.org/abs/2608.24827v2}{arXiv:2608.24827v2},
2 September 2026, Theorems 1.1--1.2 and Section 6.
\bibitem{Suzuki} M. Suzuki, \emph{Weil's quadratic form via the screw
function}, \href{https://arxiv.org/abs/2606.09096v3}{arXiv:2606.09096v3},
23 September 2026, Theorems 1.1 and 1.4.
\bibitem{CC} A. Connes and C. Consani, \emph{Weil positivity and Trace
formula, the archimedean place},
\href{https://arxiv.org/abs/2006.13771v1}{arXiv:2006.13771v1},
24 June 2020, Theorem 6.11.
\bibitem{FLINT} FLINT project, \emph{FLINT 3.6.0},
\href{https://raw.githubusercontent.com/flintlib/flint/v3.6.0/doc/source/acb_dirichlet.rst}
{versioned \texttt{acb\_dirichlet} documentation},
\texttt{zeta\_zeros}, \texttt{zeta\_nzeros}, Turing count with multiplicity.
\end{thebibliography}
\end{document}
